
\documentclass{article}
\usepackage{colortbl}
\usepackage{makecell}
\usepackage{multirow}
\usepackage{supertabular}

\begin{document}

\newcounter{utterance}

\centering \large Interaction Transcript for game `dond', experiment `semi\_en', episode 6 with Qwen3.5{-}4B.
\vspace{24pt}

{ \footnotesize  \setcounter{utterance}{1}
\setlength{\tabcolsep}{0pt}
\begin{supertabular}{c@{$\;$}|p{.15\linewidth}@{}p{.15\linewidth}p{.15\linewidth}p{.15\linewidth}p{.15\linewidth}p{.15\linewidth}}
    \# & $\;$A & \multicolumn{4}{c}{Game Master} & $\;\:$B\\
    \hline

    \theutterance \stepcounter{utterance}  
    & & \multicolumn{4}{p{0.6\linewidth}}{
        \cellcolor[rgb]{0.9,0.9,0.9}{
            \makecell[{{p{\linewidth}}}]{
                \texttt{\tiny{[P1$\langle$GM]}}
                \texttt{You are playing a negotiation game in which you have to agree on how to divide a set of items among you and another player.} \\
\\ 
\texttt{Rules:} \\
\texttt{(a) You and the other player are given a set of items. Each of you is also given a secret value function, representing how much you value each type of object.} \\
\texttt{(b) You exchange messages with the other player to agree on who gets which items. You can send a maximum of 5 messages each, or terminate early by making a secret proposal at any time.} \\
\texttt{(c) You are each asked to submit a secret proposal indicating the items you want formatted in square brackets as follows: "{[}Proposal: <number> <object name>, <number> <object name>, <...>{]}"} \\
\texttt{(d) If your proposals are complementary, i.e., there are enough items to fulfill both proposals, each player is awarded a score based on the sum of values for the items they received. Otherwise, both of you get zero points.} \\
\texttt{(e) Your goal is to maximize the score you receive.} \\
\\ 
\texttt{Let us start.} \\
\\ 
\texttt{The set of available items is:} \\
\\ 
\texttt{2 eggs, 1 picture, 2 dice} \\
\\ 
\texttt{Your secret value function is:} \\
\\ 
\texttt{egg: 1, picture: 2, dice: 3} \\
\\ 
\texttt{IMPORTANT: Your messages, unless it is the secret proposal, are directly transmitted to the other player, so do not include any response to the rules or text announcing your message. To make a secret proposal, use the indicated format. Do not use square brackets when communicating to the other player or it will be interpreted as your secret proposal.} \\
            }
        }
    }
    & \\ \\

    \theutterance \stepcounter{utterance}  
    & \multicolumn{4}{p{0.6\linewidth}}{
        \cellcolor[rgb]{0.9,0.9,0.9}{
            \makecell[{{p{\linewidth}}}]{
                \texttt{\tiny{[P1$\rangle$GM]}}
                \texttt{I will propose a division that maximizes my value while ensuring the other player also receives a positive value (to avoid the zero{-}point penalty).} \\
\\ 
\texttt{My values:} \\
\texttt{{-} 2 eggs: 2} \\
\texttt{{-} 1 picture: 2} \\
\texttt{{-} 2 dice: 6} \\
\texttt{Total: 10} \\
\\ 
\texttt{If I take the 2 dice (value 6), I leave the other player with 2 eggs and 1 picture.} \\
\texttt{The other player's values (assuming they are rational and want to maximize their own score, though I don't know their specific function, I must assume they value these items positively to accept):} \\
\texttt{{-} 2 eggs: 2} \\
\texttt{{-} 1 picture: 2} \\
\texttt{Total for them: 4.} \\
\\ 
\texttt{This split (Me: 6, Them: 4) is better for me than taking everything (10) because if I take everything, they get 0, and the rule states "If your proposals are complementary... Otherwise, both of you get zero points." Since I cannot know their exact values, I must propose a split where they are likely to accept (i.e., they get something valuable).} \\
\\ 
\texttt{Let's try to give them the items I value least, or a balanced split.} \\
\texttt{Option A: I take 2 dice (6). They take 2 eggs + 1 picture (4).} \\
\texttt{Option B: I take 1 picture + 2 eggs (4). They take 2 dice (6).} \\
\texttt{Option C: I take 2 dice + 1 picture (8). They take 2 eggs (2).} \\
\\ 
\texttt{Option A seems safest as it gives the other player a significant amount (4) without giving me the maximum possible (10), which might trigger a rejection if they value the dice highly. However, since I don't know their values, the safest bet to ensure a non{-}zero outcome is to give them the items I value least, or a fair split.} \\
\\ 
\texttt{Actually, the rule says "If your proposals are complementary... Otherwise... zero points." This implies if I propose "I want X" and they propose "I want Y", and X and Y don't cover all items or overlap in a way that leaves nothing for the other, it fails. But usually, in these games, a proposal is "I want these items". If I say "I want 2 dice", and they say "I want 2 eggs and 1 picture", that is complementary.} \\
\\ 
\texttt{Let's propose I take the 2 dice. This leaves the 2 eggs and 1 picture for them.} \\
\texttt{Message: I want the 2 dice.} \\
\\ 
\texttt{Wait, the format for the proposal is: `{[}Proposal: <number> <object name>, <number> <object name>, <...>{]}`} \\
\texttt{And I must NOT use square brackets when communicating to the other player.} \\
\texttt{So I should just say: "I want the 2 dice."} \\
\\ 
\texttt{Let} \\
            }
        }
    }
    & & \\ \\

    \theutterance \stepcounter{utterance}  
    & & & \multicolumn{2}{p{0.3\linewidth}}{
        \cellcolor[rgb]{0.9,0.9,0.9}{
            \makecell[{{p{\linewidth}}}]{
                \texttt{\tiny{[GM$|$GM]}}
                \texttt{proposal does not include number/name pairs} \\
            }
        }
    }
    & & \\ \\

\end{supertabular}
}

\end{document}
